\begin{table}[H]
\centering
\caption{\textbf{Segmentation Module.}}
\begin{tabular}{|l|p{8cm}|}
\hline
\textbf{Step} & \textbf{Purpose} \\
\hline
Validity mask computation & Computes a binary mask that identifies valid regions of the image, used throughout the pipeline to exclude corrupted or irrelevant areas and prevent them from influencing segmentation. \\
\hline
Preprocessing & Enhances contrast and standardizes intensity distributions to reduce inter-image variability; converts 16-bit images to 8-bit and inverts intensities so DNA molecules exhibit higher grayscale values than the background. \\
\hline
Bilateral Filter & Reduces background noise while preserving edge information, ensuring that thin DNA molecules remain well-defined for subsequent processing. \\
\hline
Unsharp Masking & Enhances edges and local contrast to better highlight DNA structures against the background. \\
\hline
Hysteresis Thresholding & Generates a binary segmentation by separating DNA molecules from the background. \\
\hline
Area-based Pruning (I) & Removes small, spurious segmented regions based on a predefined area threshold. \\
\hline
Morphological Closing & Fills small gaps and discontinuities between segmented regions to improve the structural continuity of DNA molecules. \\
\hline
Area-based Pruning (II) & Removes small, spurious segmented regions based on a predefined area threshold, which can differ from the one used in the first pruning stage. \\
\hline
\end{tabular}
\label{tab:segmentation_pipeline}
\end{table}